% =====================================================================
%  P2.  No weight extracts C(N): a no-go for the divisor-switch route,
%       with the exact identities of its design space.
%
%  Compiles standalone with pdflatex, run twice.
% =====================================================================
\documentclass[11pt,a4paper]{article}

\usepackage{amsmath}
\usepackage{amssymb}
\usepackage{amsthm}
\usepackage[margin=1in]{geometry}
\usepackage[colorlinks=true,linkcolor=blue,citecolor=blue]{hyperref}

\theoremstyle{plain}
\newtheorem{theorem}{Theorem}
\newtheorem{corollary}[theorem]{Corollary}
\newtheorem{proposition}[theorem]{Proposition}
\newtheorem{lemma}[theorem]{Lemma}
\theoremstyle{definition}
\newtheorem{observation}[theorem]{Observation}
\newtheorem{measurement}[theorem]{Measurement}
\theoremstyle{remark}
\newtheorem{note}[theorem]{Note}

\renewcommand{\SS}{\mathfrak{S}}
\newcommand{\AAA}{\mathfrak{A}}
\DeclareMathOperator{\rad}{rad}
\DeclareMathOperator{\lcm}{lcm}

% typographic slack: let TeX stretch a line rather than run into the margin
\emergencystretch=3em

\title{No weight that Bombieri--Vinogradov reaches extracts
$\sum_{n<N}\Lambda(n)\mu(N-n)$:\\
a no-go for the divisor-switch route,\\
with the exact identities of its design space}
\author{Youngrok Lee\\
\small Independent researcher, Republic of Korea}
\date{}

\begin{document}
\maketitle

\begin{abstract}
Let $N$ be a large even integer and
$C(N)=\sum_{n<N}\Lambda(n)\mu(N-n)$. The Huang--Li reduction of binary
Goldbach to Elliott--Halberstam \cite{HL} consumes its hypothesis
through a M\"obius-weighted correlation sum in the fixed class
$n\equiv N\ (k)$, carrying a weight $w_k$. Two weights occur: $w_k=1$,
for which the sum is unconditionally small, and $w_k=\log k$, for which
it is equivalent to Goldbach itself \cite{P1}. Is there a weight
between them --- one whose divisor sum is still cheap but whose
main-term coefficient is still of size $1$, which would therefore
extract $C(N)$ from Bombieri--Vinogradov alone?

No such weight exists among those whose transform
Bombieri--Vinogradov reaches. Writing $b=\mu*w$, the two requirements
pull apart: extraction needs mass in $b$ at $K=N^{\theta'}$, while the
residual is accessible whenever $b$ is supported below
$N^{\theta'-1/2}$, and the thresholds are separated by $N^{1/2}$.
Theorem~\ref{thm:D} turns that into
$\lVert b\rVert_1/\lvert B_w\rvert\gg\exp(c_1\sqrt{(\frac12+\delta)\log
N})$, which exceeds every power of $\log N$ and destroys the extraction
\emph{even when the hard input is granted for free}.
Theorem~\ref{thm:Dprime} shows the separation survives granting
$EH(N^{\theta_E})$ for any fixed $\theta_E<1$, so the obstruction is
not a shortage of level. The polynomial family $w_k=f(\log k)$ lies
outside the support hypothesis and is closed separately, by a different
and weaker argument.

\emph{What is not shown.} The result is for one evaluation of the
switch identity, the triangle inequality; an evaluation exploiting sign
correlations among that identity's three terms is not excluded, and
nothing here says such correlations are absent. Outside the reached
transforms the statement is weaker and is marked as such. The switch
identity itself is taken from \cite{P1} and not reproved here, so half
of Theorem~\ref{thm:D} is conditional on it.

The design space is laid out exactly first: six finite rearrangements
and one consequence of them identify the per-modulus discrepancy as a
dilate of a M\"obius--prime correlation, show the signs $\mu(k)$ cancel
out of the $\log k$ branch exactly, and exhibit the Goldbach count as a
signed sum of nonnegative layers. These are what make the result a
statement about a space rather than about two examples. All six are
formalised in Lean~4 against Mathlib.

Net progress toward Goldbach: zero. The value of a no-go of this kind
is that it closes a design space rather than leaving it to be
re-explored.
\end{abstract}

% =====================================================================
\section{Introduction}

\subsection{The setting}

Let $N$ be a large even integer, $\theta'\in(1/2,1)$ fixed, and
$K=N^{\theta'}$. Put
\begin{equation}\label{eq:C}
  C(N)=\sum_{n<N}\Lambda(n)\mu(N-n),
  \qquad
  A(N;k)=\sum_{\substack{n<N\\ n\equiv N\ (k)}}\Lambda(n)\mu(N-n),
\end{equation}
and $\tilde r(N)=\sum_{n<N}\Lambda(n)\Lambda(N-n)$ for the prime-pair
sum, so that $\tilde r(N)>0$ for all large even $N$ is binary
Goldbach. For a weight $w:\mathbb N\to\mathbb R$ put
\begin{equation}\label{eq:Tw}
  T_w \;=\;\sum_{\substack{k<K\\ (k,N)=1}}\mu(k)\,w_k
     \Bigl[A(N;k)-\frac{C(N)}{\varphi(k)}\Bigr].
\end{equation}
The functional $E_3$ of \cite{HL} is $T_w$ at $w_k=\log k$; $E_4$ is
not $T_w$ at $w_k=1$ but is obtained from it by partial summation, an
inner factor $\log(N-n)$ separating the two \cite[Corollary on the collapse]{P1}. $E_4$ is
unconditionally $\ll_A N(\log N)^{-A}$, and $E_3$ is, by an
unconditional identity, equivalent to Huang--Li's equation~(22) and
hence already gives binary Goldbach \cite{P1}. This paper closes the
part of the space between the two weights that its input reaches, and
records what is left open. What is reached is the weights whose
transform $b=\mu*w$ is supported below $N^{\theta'-1/2}$; the
polynomial family is closed separately and by a weaker argument, and
the rest --- including $w_k=\log k$ itself --- is closed nowhere here.
\S\ref{sec:notclaimed} says which is which.

Throughout, $\varphi$ is Euler's function, $\rad(u)=\prod_{p\mid u}p$,
$c,c_1,\dots$ are positive constants depending at most on the
parameters each statement fixes ($\theta'$, $\delta$, $\theta_E$) and
not on $N$, and
\begin{equation}\label{eq:AN}
  \AAA(N)=\prod_{q\,\nmid\,N}\Bigl(1-\frac{1}{q(q-1)}\Bigr),
  \qquad
  \SS(N)=2\prod_{p>2}\Bigl(1-\tfrac{1}{(p-1)^2}\Bigr)
          \prod_{\substack{p\mid N\\ p>2}}\Bigl(1+\tfrac{1}{p-2}\Bigr).
\end{equation}

\subsection{The design space}

Let $w$ be arbitrary and let $b:=\mu*w$, so that
$w_k=\sum_{d\mid k}b_d$; every weight has this form, and $b$ is
determined by $w$. Put
\begin{equation}\label{eq:Bw}
  B_w:=\sum_{\substack{k<K\\(k,N)=1}}\frac{\mu(k)w_k}{\varphi(k)},
  \qquad
  \|b\|_1:=\sum_{u<N}|b_u| .
\end{equation}
Running the divisor switch of \cite{P1} with the weight $w$ gives the
identity below. Its derivation is \cite{P1}'s and is not repeated
here; every negative conclusion of this paper is read off it, so a
reader checking this paper alone should treat it as a hypothesis taken
from outside. It is not the only one. The others, each used where it is
named, are: the bound $\rho_n(x)\ll e^{-c_1\sqrt{\log x}}$ of
Huang--Li's Lemma~1, on which Lemma~\ref{lem:extract} turns;
\cite{P1}'s residual identity and Bombieri--Vinogradov in the form of
\cite{MoV}, both inside Lemma~\ref{lem:bv}; \cite{P1}'s theorem on the
identity, used in Proposition~\ref{prop:direct}; and Vaughan's
$\|S_\Lambda\|_1\gg N^{1/2}$ \cite{Vau88}, used in
Proposition~\ref{prop:E}. What is
particular to \eqref{eq:extract} is that it is the one every negative
conclusion passes through.

Which of the four is load-bearing is a different question from which
is most used, and the answer is not \eqref{eq:extract}. If
\eqref{eq:extract} were wrong, \eqref{eq:loss} would survive --- its
proof uses Lemma~\ref{lem:extract} and not the switch identity --- and
what would die is the step from \eqref{eq:loss} to a bound on
$|C(N)|$. If Bombieri--Vinogradov reached higher moduli than assumed,
the loss factor would shrink but not vanish, and
Theorem~\ref{thm:Dprime} already measures that. The $\rho$ decay is the
one whose failure takes several statements at once:
Lemma~\ref{lem:extract} turns on it, hence \eqref{eq:loss} and with it
Theorem~\ref{thm:D}, Theorem~\ref{thm:Dprime} and
Proposition~\ref{prop:Dpp}(i); and Proposition~\ref{prop:direct}'s rate
$B_{\log}(K)+\SS(N)\ll e^{-c_1\sqrt{\log K}}\log K$ is the same bound
applied to $\rho_N$, so \eqref{eq:directcond} goes too. That is the
single most expensive dependence here, and it is not checked in this
packet --- Measurement~\ref{meas:extract} checks
Lemma~\ref{lem:extract} as an identity at $b=\delta_{d_0}$, which
tests the algebra and not the decay.
\begin{equation}\label{eq:extract}
  B_w\cdot C(N)\;=\;
  \underbrace{\sum_{u<N}\Lambda(N-u)\mu^2(u)\,b_u}_{\text{complete part}}
  \;-\;\underbrace{\mathcal R_w}_{\text{residual}}
  \;-\;\underbrace{T_w}_{\text{the object itself}}
  \;+\;O\!\left(N^{o(1)}\right),
\end{equation}
the complete part being evaluated by Lemma~\ref{lem:Gb} below, which
identifies $b$ as exactly the complete divisor transform. The three
named terms are not defined here in the same way: the complete part
and $T_w$ are written out above and below, while $\mathcal R_w$ is the
residual of \cite{P1}'s switch, and its formula is not reproduced in
this paper. What is used of it is only what
Lemma~\ref{lem:bv} states --- that expanding $w_k=\sum_{d\mid k}b_d$
splits it as $\sum_db_d\mathcal R^{(d)}$ with each
$\mathcal R^{(d)}$ a sum of $\Lambda$ over progressions to moduli $md$.
A reader wanting the residual itself has to open \cite{P1}.

\begin{note}[the third term is not a remainder]\label{note:Tw}
The term $T_w$ in \eqref{eq:extract} is \eqref{eq:Tw} itself, and it
cannot be absorbed into the $O(N^{o(1)})$: at $w=\log$ it \emph{is}
$E_3$, which
\cite[Theorem on the identity]{P1} evaluates as
$\tilde r(N)-\SS(N)(N-C(N))+O_A(N(\log N)^{-A})$. Writing
\eqref{eq:extract} without it turns
a rearrangement into a false statement --- the two sides then differ by
$0.44N$, $0.38N$, $0.32N$ at $N=2\cdot10^5,\,4\cdot10^5,\,8\cdot10^5$
with $\theta'=0.56$, which are exactly the measured values of
$|T_{\log}|/N$ at those $N$, read from the run
Measurement~\ref{meas:flatsum} cites rather than from that
measurement's own printed row; over the eight $N$ that run sweeps, the
same ratio falls to $0.1245$, so over that range the two sides differ
by a decreasing fraction of $N$ and not by a constant multiple of it.
That the decrease is a power of $N$ is a fit to those eight points and
not a bound; the use made of it here is only the lower one, that
$|T_{\log}|$ stays above $0.12N$ throughout. The $w=1$ case of the same
collection, $T_1=\sum H-C(N)B(K)$, is checked numerically and holds to
a worst relative error $1.875\cdot10^{-16}$
(\texttt{lab\_weight\_gap.py}); \eqref{eq:extract} itself is not
checked numerically anywhere here.

Keeping it explicit is what makes Theorem~\ref{thm:D} say something.
The identity by itself yields nothing about $C(N)$: an unconditional
bound on $T_w$ is exactly what the whole problem asks for, and the only
weight for which one is known is $w=1$, by \cite[Theorem on the flat branch]{P1}. The
theorem below therefore \emph{grants} the bound on $T_w$ and shows the
extraction fails anyway.
\end{note}

\begin{lemma}\label{lem:Gb}
For squarefree $u$, $\sum_{k\mid u}\mu(k)w_k=\mu(u)\,b_u$.
\end{lemma}

\begin{proof}
$\sum_{k\mid u}\mu(k)\sum_{d\mid k}b_d
 =\sum_{d\mid u}b_d\,\mu(d)\sum_{j\mid(u/d)}\mu(j)
 =\sum_{d\mid u}b_d\,\mu(d)\,[\,u/d=1\,]=\mu(u)b_u$,
using $(d,u/d)=1$ for squarefree $u$.
\end{proof}

The two known cases are the two ends of this space: $w=1$ gives
$b=\delta_1$, $\|b\|_1=1$ and (Huang--Li's Lemma~1)
$B_w\ll e^{-c\sqrt{\log K}}$; while $w=\log$ gives $b=\Lambda$,
$\|b\|_1\asymp N$ --- the complete part \emph{is} the binary Goldbach
sum --- and $B_w\to-\SS(N)\asymp1$.

\subsection{Results}

Section~\ref{sec:identities} records the exact structure of the space:
the per-modulus discrepancy is a dilate (Proposition~\ref{prop:dilate}),
the signs $\mu(k)$ cancel exactly on the $\log k$ branch
(Proposition~\ref{prop:posweights}), both known weights are the same
sum of dilated walls with different weights
(Proposition~\ref{prop:flatsum}), the untruncated object is the
Goldbach count itself (Proposition~\ref{prop:untrunc}), which is
therefore a signed sum of nonnegative layers
(Proposition~\ref{prop:layers}) over Bombieri--Vinogradov
moduli carrying no M\"obius (Proposition~\ref{prop:combined}); and
the wall cancels out of the count entirely
(Proposition~\ref{prop:direct}).

Section~\ref{sec:nogo} proves the no-go: Lemmas~\ref{lem:extract}
and~\ref{lem:bv} are the two opposing constraints,
Theorem~\ref{thm:D} is their consequence, and
Theorem~\ref{thm:Dprime} shows it survives Elliott--Halberstam.
Section~\ref{sec:Dpp} closes polynomial weights, which lie outside
the hypothesis of Theorem~\ref{thm:D}. Section~\ref{sec:circle}
records the position of the circle method.

\subsection{What is not claimed}\label{sec:notclaimed}

\begin{itemize}
\item \textbf{This is a no-go for one method, and for the part of its
weight space that method can evaluate.} Theorem~\ref{thm:D}
closes divisor switching with Bombieri--Vinogradov as the only input,
for every weight whose transform that input reaches. Of the weights
outside that hypothesis, the polynomial family $w_k=f(\log k)$ is
addressed by Proposition~\ref{prop:Dpp}, whose (ii) is a statement of
what is not known rather than an impossibility; the weights that are
neither reached by the hypothesis nor polynomial are closed nowhere
here. It is not, and does not
claim to be, an obstruction to other methods --- in particular it says
nothing about the spectral large sieve of \cite{Li23}.

\item \textbf{The genre is classical.} Bombieri's asymptotic sieve
\cite{Bom76} shows that sieve weights alone cannot detect primes,
however chosen, and the parity obstruction it isolates is the reason
one expects a statement of this shape. What is added here is the
precise quantitative form for this route
(Note~\ref{note:bombieri}).

\item \textbf{Proposition~\ref{prop:E} is about two pairings, not
about the circle method.} What is shown is that the $(2,2)$ and
$(\infty,1)$ H\"older pairings of the integral lie at or above the
trivial bound. A circle-method treatment of a binary problem does not
ordinarily proceed by a global pairing at all: it splits into major and
minor arcs and treats the minor arcs bilinearly. The Parseval floor is
taken over the whole circle and so does not depend on where that split
is drawn, but a bilinear treatment of the minor arcs is a different
mechanism and is not addressed here. Nothing below measures the major
arcs or asserts anything about their size.

\item \textbf{No progress toward Goldbach.} Nothing here bounds
$C(N)$, and nothing here is a step toward bounding it.

\item \textbf{The identities of Section~\ref{sec:identities} are
rearrangements.} Six of the seven are exact and finite, none of those
is an estimate, and none makes any quantity smaller.
Proposition~\ref{prop:direct} is the exception: it carries an
$O_A(N(\log N)^{-A})$ and uses an asymptotic bound on $|C(N)|$, so it
is a consequence of the rearrangements rather than one of them.
\end{itemize}

% =====================================================================
\section{The design space, exactly}\label{sec:identities}

\begin{proposition}[the discrepancy is a dilate]\label{prop:dilate}
Unconditionally,
\begin{equation}\label{eq:dilate}
  A(N;k)\;=\;\mu(k)\,H(N;k),
  \qquad
  H(N;k)\;=\;\sum_{\substack{m<N/k\\ (m,k)=1}}\Lambda(N-mk)\,\mu(m).
\end{equation}
\end{proposition}

\begin{proof}
The class $n\equiv N\pmod k$ is $k\mid N-n$, so writing $N-n=mk$ turns
the sum into $\sum_{m<N/k}\Lambda(N-mk)\mu(mk)$; and
$\mu(mk)=\mu(m)\mu(k)$ when $(m,k)=1$, zero otherwise.
\end{proof}

Identity~\eqref{eq:dilate} qualifies the shape of the whole problem.
$EH_\mu$ is consumed where one M\"obius argument \emph{divides} the
other, but the per-modulus discrepancy is $\mu(k)$ times a sum in
which the two arguments are related by a \emph{difference}, at scale
$N/k$. The two couplings are not two fields; they are the same field
seen before and after the switch. It also explains why $\mu$ is the
noisier of the two inputs in progressions: for $\mu$,
$|A(N;k)|=|H(N;k)|$ is a M\"obius--prime correlation of length $N/k$
and nothing makes it small, whereas for a random sign pattern
$\varepsilon$ the sum $\sum_m\Lambda(N-mk)\varepsilon(mk)$ is a sum of
independent signs and exhibits square-root cancellation.

\begin{proposition}[the weights are nonnegative]\label{prop:posweights}
Substituting \eqref{eq:dilate} into $T_{\log}$ and using $\mu(k)^2=1$
on squarefree $k$,
\begin{equation}\label{eq:posweights}
  T_{\log}\;=\;\sum_{\substack{k<K\\(k,N)=1}}
    \mu^2(k)\,(\log k)\,H(N;k)\;-\;C(N)\,B_{\log}(K),
  \qquad
  B_{\log}(K)=\!\!\sum_{\substack{k<K\\(k,N)=1}}\!\!
     \frac{\mu(k)\log k}{\varphi(k)} .
\end{equation}
\end{proposition}

\begin{proof}
Immediate from \eqref{eq:dilate} and $\mu(k)\mu(k)=\mu^2(k)$.
\end{proof}

The signs $\mu(k)$ cancel exactly: they cancel against the $\mu(k)$
that \eqref{eq:dilate} extracts from the progression sum, and what is
left is a sum of dilated M\"obius--prime correlations with
\emph{nonnegative} weights $\mu^2(k)\log k$. Huang--Li discard the
$\mu(k)$ by the triangle inequality; on the $w_k=1$ branch keeping
them is what makes the unconditional theorem of \cite{P1} possible; on
the $\log k$ branch they are not there to keep. Whatever smallness
$T_{\log}$ has must therefore come from the signs of $H(N;k)$ across
$k$. It does not come from there, and the amount is measured. Writing
$G=\sum_k(\log k)|H|\big/\bigl|\sum_k(\log k)H\bigr|$ for the gain that
cancellation across dilations buys, $G$ reads $1.834$, $1.804$,
$2.207$, $2.588$ and $2.789$ over $N=2\cdot10^5$ to $3.2\cdot10^6$,
where the sums run over $313$ to $1485$ moduli
(\texttt{lab\_positive\_weights.py}).

\emph{What independent signs would give is not $\sqrt{\#k}$.} Writing
$x_k=(\log k)H(N;k)$, the quantity $G$ is $\lVert x\rVert_1/
\lvert\sum_kx_k\rvert$, and under independent signs
$\lvert\sum_kx_k\rvert$ is of the size of $\lVert x\rVert_2$; so the
reference for $G$ is $\lVert x\rVert_1/\lVert x\rVert_2$. That equals
$\sqrt{\#k}$ only when every $\lvert x_k\rvert$ is the same, and
Cauchy--Schwarz makes it at most $\sqrt{\#k}$ otherwise. Here the
magnitudes are not equal: the same run prints
$\lVert x\rVert_1/\lVert x\rVert_2$ as $11.96$, $14.23$, $17.90$,
$21.43$ and $26.21$, which is $0.66$ to $0.69$ of the count-based
figure. \textbf{A count of moduli overstates the reference by about a
half.}

Against the reference that is correct for these magnitudes, the
measured $G$ of $1.8$ to $2.8$ is still smaller by an order of
magnitude, so the conclusion is unchanged and only the comparison
moves. The step from a signed sum to its absolute counterpart --- which
is what \eqref{eq:directcond} makes below --- costs a factor near two
on this range, where independence would have bought between twelve and
twenty-six.

\begin{proposition}[both ends are the same sum of dilated walls]%
\label{prop:flatsum}
For any weight $w$,
\begin{equation}\label{eq:flatsum}
  T_w\;=\;\sum_{\substack{k<K\\(k,N)=1}}\mu^2(k)\,w_k\,H(N;k)
   \;-\;B_w\,C(N).
\end{equation}
In particular
\[
  T_1=\!\!\sum_{\substack{k<K\\(k,N)=1}}\!\!\mu^2(k)H(N;k)-B_1C(N),
  \qquad
  T_{\log}=\!\!\sum_{\substack{k<K\\(k,N)=1}}\!\!\mu^2(k)(\log k)H(N;k)
   -B_{\log}C(N),
\]
and both weights are nonnegative on the range.
\end{proposition}

\begin{proof}
Put \eqref{eq:dilate} into \eqref{eq:Tw}. The transform carries
$\mu(k)$, so terms with $\mu(k)=0$ contribute nothing on the left; on
the remaining $k$, $\mu(k)A(N;k)=\mu(k)^2H(N;k)$. The restriction is
therefore carried on the right, and the subtracted mean term assembles
into $B_wC(N)$ by \eqref{eq:Bw}.
\end{proof}

\begin{note}[the restriction is not decoration]
Dropping $\mu^2(k)$ from \eqref{eq:flatsum} does not weaken it, it
falsifies it: $H(N;k)$ does not vanish at squarefull $k$, while the
left-hand side does. The terms so admitted are of the size of $T_w$
itself --- at $\theta'=0.56$ and $N=2\cdot10^{5}$ they contribute
$-0.054316\,N$ at $w=1$ and $-0.287165\,N$ at $w=\log$, the latter
being $0.656087$ of $|T_{\log}|$ in absolute value. The fraction is not stable: over
$N=2,4,8\cdot10^{5}$ it reads $0.656087$, $0.701640$, $0.788480$, so
it rises with $N$ and the two thirds is the value at the bottom of
that range. It stays above $1/2$ throughout, which is the whole of
what this note asserts --- the restriction is load-bearing and not
cosmetic (\texttt{audit\_squarefull\_admitted.py}).
\end{note}

Since $B_1\ll e^{-c\sqrt{\log K}}$ kills the second term while
$B_{\log}\to-\SS(N)$ does not, the unconditional theorem of \cite{P1}
says exactly that the \emph{flat} sum of dilated walls is
$\ll_AN(\log N)^{-A}$, and the Goldbach problem says the same sum
weighted by $\log k$ is not small. Everything between what is proved
and what is open is the factor $\log k$ inside a positively weighted
sum of the same terms.

\begin{measurement}[the two ends, measured]\label{meas:flatsum}
Over $N=2\cdot10^5$ to $2.56\cdot10^7$ by doubling, at
$\theta'=0.56$, \eqref{eq:flatsum} holds to a worst relative error of
$1.875\cdot10^{-16}$, and
\[
  \frac{\bigl|\sum_k\mu^2(k)H(N;k)\bigr|}{\bigl|\sum_k\mu^2(k)(\log k)H(N;k)\bigr|}
  \;=\;0.1785,\ 0.1559,\ 0.1484,\ 0.1483,\ 0.1367,\ 0.1334,\
       0.1241,\ 0.1135,
\]
falling on average; at the top $|T_1|/N=0.01425$ against
$|T_{\log}|/N=0.1245$. Fitted against $\log N$ the flat sum decays as
$N^{-0.3505}$ and $|T_{\log}|/N$ as $N^{-0.2658}$: the gap widens.
Every fitted exponent in this paper is one arithmetic family's. These
$N$ are $2\cdot10^5\cdot2^{\,j}$ with $2\cdot10^5=2^6\cdot5^5$, so each
has odd radical $5$, and the same doubling carries the $-0.3620$ below
and the residual exponent of Measurement~\ref{meas:direct};
\texttt{c02\_flatsum\_k1.txt} states the field for that reason. What
these exponents do across radicals is not measured, and no step of the
no-go rests on their values. (The
neighbouring $-0.2696$ printed in the same run is the exponent of
$|\sum_k(\log k)H(N;k)|/N$, which is a different sum --- the two differ
by $C(N)B_{\log}(K)$ --- and not a second measurement of the same one.
Both fall out of one sieve and one index set, so what separates them is
the quantity and not the implementation.) The index set here runs from $k=1$, as \eqref{eq:Tw} and
\eqref{eq:flatsum} require; note that $H(N;1)=C(N)$ exactly, so
dropping that one term removes $C(N)$ from the flat sum --- which
changes the ratios above by up to $11.6\%$ --- taken over the $k\ge1$
column, which is the one those ratios are formed from; against the
$k\ge2$ column it is $10.4\%$. Those two percentages are not printed by
any run: they are the largest relative gap between the
\texttt{ratio\_code} and \texttt{ratio\_paper} columns of
\texttt{c02\_flatsum\_k1.txt}, attained at $N=4\cdot10^5$ where the
columns read $0.1740$ and $0.1559$, and each is that gap divided by one
column in turn. The fitted exponent moves to
$-0.3620$, while leaving $|T_1|$ untouched to $10^{-17}$, since the
$k=1$ term of $T_1$ is $A(N;1)-C(N)/\varphi(1)=0$. The two columns are
computed side by side, from a sieve built independently of the one
\texttt{lab\_weight\_gap.py} uses, in
\texttt{c02\_flatsum\_k1.py}; that script's loop is the $k\ge1$ one,
and \texttt{lab\_weight\_gap.py}'s starts at $k=2$. Identity~\eqref{eq:dilate} itself is verified over every
squarefree $k<N^{0.56}$ coprime to $N$ at
$N=2\cdot10^5,\,4\cdot10^5,\,8\cdot10^5$ with worst relative error
$2.294\cdot10^{-14}$, $1.197\cdot10^{-13}$, $4.591\cdot10^{-14}$, and
\eqref{eq:posweights} against a direct evaluation of $T_{\log}$ to at
most $1.8\cdot10^{-16}$ over $N=2\cdot10^5$ to $3.2\cdot10^6$
(\texttt{lab\_weight\_gap.py}, \texttt{lab\_dilate\_identity.py},
\texttt{lab\_positive\_weights.py}).
\end{measurement}

\begin{proposition}[the untruncated sum is the count]%
\label{prop:untrunc}
Summing over \emph{all} $k$, with no truncation and no coprimality
restriction,
\begin{equation}\label{eq:untrunc}
  \sum_{k\ge1}(\log k)\,\mu(k)\,A(N;k)
   \;=\;-\sum_{n<N}\Lambda(n)\,\mu(N-n)\Lambda(N-n)
   \;=\;\sum_{p<N}\Lambda(N-p)\log p .
\end{equation}
So the untruncated object \emph{is} $\tilde r(N)$, up to its
prime-power part.
\end{proposition}

\begin{proof}
Both exchanges are finite sums. Since $A(N;k)$ sums $n$ over
$k\mid N-n$, exchanging the order and applying
$\sum_{d\mid u}\mu(d)\log d=-\Lambda(u)$ --- M\"obius inversion of
$\log=\mathbf 1*\Lambda$ --- gives the first equality. The second
holds because $\mu(u)\Lambda(u)=-\log p$ at $u=p$ prime and vanishes
at every higher prime power.
\end{proof}

Proposition~\ref{prop:untrunc} is not a decoration: it says that
\eqref{eq:flatsum} with $w=\log$ is a \emph{rearrangement} of
$\tilde r(N)$ rather than an approximation waiting for an estimate,
so all of its content is in the truncation at $K$.

\begin{proposition}[the count as signed layers]\label{prop:layers}
Cutting the same double sum along the inner variable of
\eqref{eq:dilate},
\begin{equation}\label{eq:layers}
  \sum_{k}\mu^2(k)(\log k)H(N;k)\;=\;\sum_{m}\mu(m)\,L(N;m),
  \qquad
  L(N;m):=\!\!\sum_{\substack{k<N/m,\ (k,m)=1\\ \mu(k)\neq0}}\!\!
            (\log k)\,\Lambda(N-mk),
\end{equation}
and every $L(N;m)$ is nonnegative. With
Proposition~\ref{prop:untrunc}, the Goldbach count is therefore an
alternating sum, signed by $\mu$, of nonnegative layers.
\end{proposition}

\begin{proof}
A finite exchange; the coprimality conditions on the pair $(k,m)$ are
symmetric. Nonnegativity is that of $\log k$ on $k\ge1$ and of
$\Lambda$. The exchange and the nonnegativity are checked numerically
in \texttt{lab\_layer\_decomposition.py}, which also measures how much
of the total the layers below the truncation carry.
\end{proof}

\begin{proposition}[the count over a combined modulus]%
\label{prop:combined}
Expanding the squarefree condition inside \eqref{eq:layers} by
$\mu^2(k)=\sum_{d^2\mid k}\mu(d)$ and writing $k=d^2j$,
\begin{equation}\label{eq:combined}
  \sum_{p<N}\Lambda(N-p)\log p
   \;=\!\!\sum_{\substack{m,d\ \text{squarefree}\\ (d,m)=1,\ md^2<N}}\!\!
     \mu(m)\mu(d)\,L_2(N;m,d),
\end{equation}
where
\[
  L_2(N;m,d)=\!\!\sum_{\substack{j<N/(md^2)\\ (j,m)=1}}\!\!
   \log(d^2j)\,\Lambda\bigl(N-md^2j\bigr),
\]
and $L_2$ carries no squarefree condition. It is not quite a single
progression: the surviving $(j,m)=1$ unfolds by M\"obius, $j=ej'$ with
$e\mid m$, into progressions to moduli $md^2e$. Only squarefree $e$
survive $\mu(e)$, so $e\mid\rad(m)$ and $L_2$ is a signed sum of at
most $2^{\omega(m)}$ Bombieri--Vinogradov objects with moduli running
up to $md^2\rad(m)$ --- equal to $m^2d^2$ exactly when $m$ is
squarefree. What the squarefree condition bought is that none of them
carries a $\mu$; what it did not buy is a single modulus.
\end{proposition}

\begin{proof}
Finite exchange, and $(k,m)=1$ splits as $(d,m)=(j,m)=1$.
\end{proof}

\begin{measurement}[how large that modulus has to be]\label{meas:combmod}
Being free of M\"obius is a statement about shape, not about level,
and the level is the whole question. The truncation below is taken at
$md^2<Q$, while Proposition~\ref{prop:combined} unfolds that into
moduli up to $md^2\rad(m)$. So the $Q^*$ reported here is a lower bound
on the level actually needed --- the conservative direction, since what
is claimed is that the level must be \emph{large}, and a lower bound
that is already large settles that a fortiori. Verified
over $N=2\cdot10^5$ to $1.6\cdot10^6$ --- $166384$ to $1331048$ pairs
--- \eqref{eq:combined} holds to a worst relative error of
$4.525\cdot10^{-15}$. Truncating it to $q<Q$ and letting $Q^*(N)$ be
the least $Q$ past which the deficit stays under $\varepsilon N$ at the
chosen tolerance $\varepsilon=0.01$,
\[
  Q^*=160112,\ 200848,\ 482589,\ 589203,
  \qquad
  Q^*/\sqrt N=358.0213,\ 317.5686,\ 539.5509,\ 465.8059,
\]
so at that tolerance $Q^*$ exceeds $\sqrt N$ by two orders at every $N$
tested. Two limits on that reading, with different grounds, and the
difference between them matters.

The first is visible in the four numbers themselves: the ratio is not
monotone in $N$: the second entry is below the first and the fourth is
below the third, while the third is above both of them. There is no
trend in $N$ here to read, at this tolerance, and none is claimed.

The second is not visible here and is not measured here. The tolerance
is chosen and not derived, and \emph{no exponent may be fitted to
$Q^*$}: the deficit oscillates while decaying, so $Q^*(\varepsilon)$
reads the last excursion above $\varepsilon N$ rather than a trend, and
the ratio has no consistent trend in $N$ across tolerances. The sweep
establishing that is not in this packet, so the statement about
tolerances is reported from the source and not established here.

What stands is therefore a comparison at one fixed tolerance and no
more: expanding the squarefree condition buys the M\"obius-free moduli and
costs two orders in the size of $Q^*$ at $\varepsilon=0.01$. No rate is
claimed, and this statistic does not measure one
(\texttt{lab\_combined\_modulus.py}).
\end{measurement}

\begin{proposition}[the wall cancels out of the count]%
\label{prop:direct}
Substituting \eqref{eq:posweights} into the identity of \cite[Theorem on the identity]{P1},
the two occurrences of $C(N)$ carry opposite signs and cancel:
\begin{equation}\label{eq:direct}
  \tilde r(N)\;=\;\SS(N)N
   \;+\!\!\sum_{\substack{k<K\\(k,N)=1}}\!\!\mu^2(k)(\log k)H(N;k)
   \;-\;C(N)\bigl(B_{\log}(K)+\SS(N)\bigr)
   \;+\;O_A\!\left(\frac{N}{(\log N)^{A}}\right).
\end{equation}
Since $B_{\log}(K)\to-\SS(N)$, the quantity $C(N)$ enters the Goldbach
count only through the vanishing factor $B_{\log}(K)+\SS(N)$.
Consequently
\begin{equation}\label{eq:directcond}
  \sum_{\substack{k<K\\(k,N)=1}}\mu^2(k)(\log k)\bigl|H(N;k)\bigr|
   \;\le\;(1-\varepsilon)\,\SS(N)\,N
\end{equation}
suffices for $\tilde r(N)>0$, for all large even $N$.
\end{proposition}

\begin{proof}
The limit $B_{\log}(K)\to-\SS(N)$ is the one thing here that
\eqref{eq:direct} does not supply, so it is proved. With
$F(s)=\prod_{p\nmid N}\bigl(1-\tfrac{p^{-s}}{p-1}\bigr)$ one has
$B_{\log}(\infty)=-F'(0)$, and $F(s)=G(s)\zeta(s+1)^{-1}$ with
$G(0)=\SS(N)$, so $F'(0)=G(0)=\SS(N)$ and the limit is $-\SS(N)$. The
rate is Huang--Li's Lemma~1 in the same form used above, applied to
$\rho_N$ and summed by parts against $\log k$ over $k\ge K$, the
partial summation being where the extra factor comes from:
$B_{\log}(K)+\SS(N)\ll e^{-c_1\sqrt{\log K}}\log K$, which is $o(1)$
uniformly in $N$ once $K=N^{\theta'}$ with $\theta'$ fixed.

Given that, the claim is immediate from \eqref{eq:direct} and
$|C(N)|\le\psi(N)\ll N$, the trivial bound, which is all that is needed
because the factor multiplying it is the $o(1)$ just bounded rather
than $\SS(N)$.
\end{proof}

\begin{note}[what the cancellation buys, and what it does not]
Condition~\eqref{eq:directcond} has the same shape as the
constant-factor condition of \cite[Proposition on the constant-factor demand]{P1}: a bound at level
$N^{\theta'}$, with no saving in $\log N$. Nothing has moved off the
level axis. What has changed is the threshold constant, from
$\SS(N)(1-\AAA(N))$ to $\SS(N)$, and the change removes the arithmetic
sensitivity along with the slack, because the two were the same thing:
$1-\AAA(N)$ collapses at primorials, and $\SS(N)$ does not. This is
offered as a weaker sufficient condition, proved from two identities.
It is not offered as something computation confirms: at accessible $N$
the unspecified $O_A$ term of \eqref{eq:direct} is numerically the
very quantity \eqref{eq:directcond} must bound.
\end{note}

\begin{measurement}[the two thresholds, compared]\label{meas:direct}
At $\theta'=0.56$ over $N=2\cdot10^5$ to $3.2\cdot10^6$, with
$B_H(N)=\sum_{k<K}\mu^2(k)(\log k)|H(N;k)|$, the new ratio
$B_H(N)/(\SS(N)N)$ reads
$0.4578,\,0.4064,\,0.4079,\,0.3769,\,0.3338$ --- under $1$ throughout,
and falling overall though not at every step, the third entry sitting
above the second --- against
$2.1519,\,1.9105,\,1.9175,\,1.7720,\,1.5691$ for the same numerator
against the old threshold, $B_H(N)/(\SS(N)(1-\AAA(N))N)$, which is
above $1$ throughout. \emph{The numerator is held fixed across the two
lists}; only the threshold constant moves, which is the whole of the
comparison. (For the numerator of \cite[the constant-factor inequality]{P1}, which keeps the
subtracted mean inside the absolute value and so is a different sum,
the old ratio reads $2.1591,\,1.9747,\,1.9500,\,1.7483,\,1.5798$ ---
the same conclusion, a different quantity.) Across a
seven-point family of arithmetically diverse $N$ the new ratio spans a
factor $2.00$, where the same family taken against the old threshold
with \emph{that} other numerator spans $23.25$. The two spans are not
formed from the same sum, and the run prints them under column names
that say so; what the spans compare is the two thresholds, and the
arithmetic sensitivity they differ by belongs to $1-\AAA(N)$, which is
the only factor between them. The residual coupling is
negligible at these $N$: $|B_{\log}+\SS|/\SS$ stays under $0.0383$ and
$|C(N)(B_{\log}+\SS)|/N$ never exceeds $3.315\cdot10^{-4}$, four
orders below $\SS(N)$. The residual of \eqref{eq:direct} decays as
$N^{-0.2794}$ with correlation $-0.99930$ and falls to $2\%$ of $N$
only near $N=10^{10.16}$
(\texttt{lab\_direct\_route.py}, \texttt{lab\_direct\_identity.py};
the two lists held to one numerator, and the third quantity in the
parenthesis, are recomputed side by side in
\texttt{a1\_bh\_vs\_b.py} --- \texttt{lab\_direct\_route.py} prints
the parenthetical list and the first of the two above it; only the
second, $2.1519,\dots$, is exclusive to \texttt{a1\_bh\_vs\_b.py}).
\end{measurement}

% =====================================================================
\section{The no-go}\label{sec:nogo}

\subsection{The two opposing constraints}

\begin{lemma}[extraction]\label{lem:extract}
\[
  B_w=\sum_{\substack{d<K\\(d,N)=1}}b_d\,
   \frac{\mu(d)}{\varphi(d)}\,\rho_{dN}\!\left(\frac{K}{d}\right),
  \qquad
  \rho_{n}(x)=\!\!\sum_{j<x,\,(j,n)=1}\!\!\frac{\mu(j)}{\varphi(j)},
\]
and $\rho_n(x)\ll e^{-c_1\sqrt{\log x}}$ whenever $\log n\ll\log x$.
\end{lemma}

\begin{proof}
Write $k=dj$; for squarefree $k$ one has $(d,j)=1$,
$\mu(k)=\mu(d)\mu(j)$ and $\varphi(k)=\varphi(d)\varphi(j)$. The bound
on $\rho$ is Huang--Li's Lemma~1, a form of
Goldston--Y{\i}ld{\i}r{\i}m \cite{GY}, and the hypothesis
$\log n\ll\log x$ is theirs.
\end{proof}

\begin{note}[the hypothesis is load-bearing, and it is met where it is used]
Without $\log n\ll\log x$ the bound is false: at $\lfloor K/d_0\rfloor=1$
the run behind Measurement~\ref{meas:extract} prints
$|B_w|\varphi(d_0)=1.000000$ over all $5130$ admissible $d_0$
(\texttt{audit\_extraction\_tradeoff.py}; that row is not among the
two the measurement reproduces), no
saving at all, and that is exactly the corner where $\log n\asymp\log N$
while $\log x=O(1)$. Theorem~\ref{thm:D} is unaffected: it maximises
only over $d\le N^{\theta'-1/2-\delta}$, where $x=K/d\ge N^{1/2+\delta}$
and $n=dN\le N^{3/2}$, so $\log n\asymp\log x$ holds throughout the
range it uses. The same applies to Theorem~\ref{thm:Dprime}.
\end{note}

\begin{lemma}[BV-accessibility]\label{lem:bv}
Expanding $w_k=\sum_{d\mid k}b_d$ inside the residual gives
$\mathcal R_w=\sum_db_d\,\mathcal R^{(d)}$, where $\mathcal R^{(d)}$
is a sum of $\Lambda$ over progressions to moduli $md$ with
$m<N^{1-\theta'}$. Bombieri--Vinogradov bounds each
$\mathcal R^{(d)}\ll_AN(\log N)^{-A}$ provided
$d\le N^{1/2-\delta}/N^{1-\theta'}=N^{\theta'-1/2-\delta}$; hence, on
that support, $\mathcal R_w\ll_A\|b\|_1\,N(\log N)^{-A}$.
\end{lemma}

\begin{proof}
The residual of the switch is \cite[the residual identity]{P1}, which after unfolding
$\mu^2$ and the coprimality condition, every term is a sum of
$\Lambda$ over a progression to modulus $q=m\lcm(d'^2,e')$ with
$m<N^{1-\theta'}$ and auxiliary divisors $d',e'$ at most a fixed power
of $\log N$ --- written with primes to keep them apart from the $d$
that indexes $b_d$;
Bombieri--Vinogradov is applied in the form of \cite{MoV}.
Inserting the extra divisor $d\mid k$ multiplies the modulus by $d$,
so the level condition $q\le N^{1/2-\delta}$ reads
$d\le N^{1/2-\delta}/\bigl(m\lcm(d'^2,e')\bigr)$; since
$\lcm(d'^2,e')\ll(\log N)^{O(1)}$, that factor is absorbed into any
$\delta>0$ and the stated bound on $d$ follows. Summing the $\ll_AN(\log N)^{-A}$ over the support of
$b$ with the trivial weight $|b_d|$ gives the claim.
\end{proof}

The two lemmas pull in opposite directions. Lemma~\ref{lem:extract}
says $B_w$ is controlled by the part of $b$ sitting \emph{at} the
truncation point $K$: mass at $d\le K^{1-\varepsilon}$ is damped by
$e^{-c\sqrt{\varepsilon\log K}}$. Lemma~\ref{lem:bv} says the residual is
reachable when $b$ lives \emph{below} $N^{\theta'-1/2}$; it does not
say that a $b$ living higher puts the residual out of reach, and
nothing here does. The two thresholds are separated by a factor
$N^{1/2}$.

\subsection{The theorem}

\begin{theorem}[no weight this input reaches extracts
$C(N)$]\label{thm:D}
Fix $\theta'\in(1/2,1)$ and $0<\delta<\theta'-\tfrac12$, and let $w$
be any weight
whose transform $b=\mu*w$ is supported in
$[1,\,N^{\theta'-1/2-\delta}]$ --- the range in which the residual is
accessible to Bombieri--Vinogradov. (For $\delta\ge\theta'-\tfrac12$
that range admits at most $b=b_1\delta_1$ --- only at
$\delta=\theta'-\tfrac12$, and nothing above it --- and the statement, though true,
is nearly empty.) Then
\begin{equation}\label{eq:loss}
  \frac{\|b\|_1}{|B_w|}\;\gg\;
   \exp\!\left(c_1\sqrt{(\tfrac12+\delta)\log N}\right).
\end{equation}
Consequently, \emph{even granted} $T_w\ll_A N(\log N)^{-A}$ for every
$A>0$ --- which is what \cite[Theorem on the flat branch]{P1} supplies unconditionally at
$w=1$, and the most that $EH_\mu$ would supply for a general $w$ ---
identity \eqref{eq:extract} yields no bound on $|C(N)|$ better than
\[
  |C(N)|\;\le\;\frac{\|b\|_1}{|B_w|}
     \cdot O_A\!\left(\frac{N}{(\log N)^{A}}\right),
\]
whose prefactor \eqref{eq:loss} bounds below by
$\exp(c_1\sqrt{(\frac12+\delta)\log N})$, which exceeds every power of
$\log N$. So the route delivers no saving of any power of
$\log N$: no weight satisfying the hypothesis above extracts
$C(N)$ by divisor switching, not even when the hard input is given
for free.

Two limits on that, both in the statement because they belong there.
Weights outside the hypothesis are not covered; Proposition~\ref{prop:Dpp}
treats the polynomial ones, and by a weaker argument. And what is
closed is the absolute-value route, at both of the two places it is
taken: $\|b\|_1$ enters \eqref{eq:loss} by bounding
$\sum_d|b_d|/\varphi(d)\,|\rho_{dN}|$ termwise, and the three terms on
the right of \eqref{eq:extract} are then bounded separately. Neither
step is a property of the identity; both are the triangle inequality,
and $\|b\|_1$ is the budget that inequality spends. An evaluation that
uses sign correlations --- among the $b_d$ in the first step, or among
the complete term, the residual and $T_w$ in the second --- does not
factor through either, and is not excluded here. Nothing in this paper
says such correlations are absent.
\end{theorem}

\begin{proof}
By Lemma~\ref{lem:extract} and $\varphi(d)\ge1$,
\[
  |B_w|\le\sum_{d\le D}\frac{|b_d|}{\varphi(d)}
    \left|\rho_{dN}\!\left(\frac Kd\right)\right|
  \le\|b\|_1\max_{d\le D}\left|\rho_{dN}\!\left(\frac Kd\right)\right|
  \ll\|b\|_1\,e^{-c_1\sqrt{\log(K/D)}},
\]
with $D=N^{\theta'-1/2-\delta}$, so that $K/D=N^{1/2+\delta}$. This is
\eqref{eq:loss}.

For the second assertion, $\|b\|_1$ and $|B_w|$ are each homogeneous of
degree one in $w$, so their ratio is scale invariant and one may
normalise $\|b\|_1=1$; that invariance is what makes the normalisation
free. It is the scale at which the two thresholds are compared. Solve
\eqref{eq:extract} for $C(N)$ and bound the three terms on the right.
The complete part is $\ll\log N$ by Lemma~\ref{lem:Gb} and
$\Lambda\ll\log N$; the residual is $\ll_AN(\log N)^{-A}$ by
Lemma~\ref{lem:bv}, the support hypothesis on $b$ being exactly what
that lemma needs; and $T_w\ll_AN(\log N)^{-A}$ is the granted
hypothesis. The last two dominate the first, so
\[
  |C(N)|\;\le\;\frac{1}{|B_w|}
     \cdot O_A\!\left(\frac{N}{(\log N)^{A}}\right)
  \;=\;\frac{\|b\|_1}{|B_w|}
     \cdot O_A\!\left(\frac{N}{(\log N)^{A}}\right).
\]
This is the bound the route produces, and it is the displayed one.
What \eqref{eq:loss} adds is a lower bound on its prefactor: the
prefactor is $\gg\exp(c_1\sqrt{(\frac12+\delta)\log N})$, which
exceeds every power of $\log N$, so no choice of $w$ in the design
space turns the displayed bound into a saving of a power of $\log N$.
The inequality is not run in the other direction --- \eqref{eq:loss}
bounds the prefactor below and so cannot be substituted into an upper
bound for $|C(N)|$; what it establishes is that the one evaluation this
route offers --- the triangle inequality applied to \eqref{eq:extract}
with the bound on $T_w$ granted --- gives nothing better than the
stated ratio. It says nothing about evaluations that do not factor
through that inequality.
\end{proof}

\begin{note}[why granting $T_w$ is the right form]\label{note:granting}
Without the granted hypothesis the theorem would be vacuous rather than
strong: \eqref{eq:extract} could then be satisfied by $T_w$ alone being
large, and it says nothing about $C(N)$ at all. Granting it isolates
what the design space actually fails at. The failure is not a shortage
of input --- the input is handed over --- but the arithmetic of two
thresholds that lie a factor $N^{1/2}$ apart. That is why raising the
level does not help either (Theorem~\ref{thm:Dprime}).
\end{note}

The loss factor is $\exp(c_1\sqrt{(\frac12+\delta)\log N})$, and the $1/2$ in the
exponent is literally the exponent of the $\sqrt N$ barrier: it is the
gap between where a weight must live to see $C(N)$ (at $N^{\theta'}$,
the truncation point) and where it may live for
Bombieri--Vinogradov to close its residual (below $N^{\theta'-1/2}$).
The two known weights exhibit the two endpoints of this space;
Theorem~\ref{thm:D} shows the interior is empty where its
hypothesis reaches, and shows why.

\begin{note}[which branch the no-go covers]\label{note:nogoscope}
The two endpoints are not on the same side of Theorem~\ref{thm:D}'s
support hypothesis, and the asymmetry is worth stating because it runs
opposite to the way the two branches are usually named. That hypothesis
asks the transform $b=\mu*w$ to be supported in
$[1,N^{\theta'-1/2-\delta}]$, and
\[
  w_k=1\ \Longrightarrow\ b=\mu*1=\delta,
  \qquad
  w_k=\log k\ \Longrightarrow\ b=\mu*\log=\Lambda ,
\]
verified elementwise to $d\le4000$. The delta is inside the hypothesis
for any positive bound; the support of $\Lambda$ runs to the truncation
point, measured at $1369$ and $2971$ against bounds $2.2$ and $2.4$ at
$\theta'=0.56$. Those bounds are $N^{\theta'-1/2}$, the $\delta=0$
endpoint, which is the largest the hypothesis ever allows; any
$\delta>0$ makes them smaller and the gap wider, so the comparison is
stated in the direction least favourable to it.
So Theorem~\ref{thm:D} covers the \emph{flat} branch
and not the $\log$ one --- the flat branch being the one
\cite[Theorem on the flat branch]{P1} settles unconditionally, and the $\log$ branch the
one left open. That is consistent rather than circular: what is
forbidden here is \emph{extracting} $C(N)$ from a weighting, and
\cite[Theorem on the flat branch]{P1} does not extract $C(N)$; it shows the weighting is
small, which is a different thing.

Below the threshold the statement is empty rather than false. For
$\theta'<1/2$ the support bound $N^{\theta'-1/2-\delta}$ is under $1$
--- the run prints $0.5$ and $0.4$ at $\theta'=0.44$ for $\delta=0$,
which is the largest the bound can be --- so no weight meets
the hypothesis at all, the delta included. That Theorem~\ref{thm:D}
says nothing there is not a gap in it: it is the same
Bombieri--Vinogradov reach that governs $\theta'>1/2$ throughout,
appearing here in the theorem's own hypothesis
(\texttt{audit\_nogo\_scope.py}).
\end{note}

\begin{note}[relation to Bombieri's asymptotic sieve]\label{note:bombieri}
The genre of Theorem~\ref{thm:D} is classical. Bombieri's asymptotic
sieve \cite{Bom76} shows that sieve weights alone cannot detect
primes, however the weights are chosen, and the parity obstruction it
isolates is the reason one expects a statement of this shape. What
Theorem~\ref{thm:D} adds is not the phenomenon but its precise form
here: the obstruction is quantified as a separation of $N^{1/2}$
between two thresholds, it is specific to the divisor-switch route,
and it survives granting the full Elliott--Halberstam conjecture
rather than merely Bombieri--Vinogradov. The reason the statement is
worth making is narrow but real: without it, the $\log k$ weight looks
like one member of a family of candidate consumptions, and the
identity of \cite[Theorem on the identity]{P1} then looks like an accident of one
choice of weight rather than a property of the route.
\end{note}

\subsection{The no-go survives Elliott--Halberstam}

The obvious objection to Theorem~\ref{thm:D} is that
Bombieri--Vinogradov is not the only available input: Huang--Li's
Theorem~1 itself \emph{assumes} $EH$ for $\Lambda$ at level
$N^{\theta}$. Raising the assumed level does not help.

\begin{theorem}\label{thm:Dprime}
Suppose $\Lambda$ has level of distribution $\theta_E\in(0,1)$, i.e.
$EH(N^{\theta_E})$ holds. Then the residual is accessible only for $b$
supported on $d\le N^{\theta_E-(1-\theta')}$, while extraction still
requires mass at $d\asymp K=N^{\theta'}$; the two thresholds are
separated by $N^{1-\theta_E}$, and
\[
  \frac{\|b\|_1}{|B_w|}\;\gg\;
  \exp\!\left(c_1\sqrt{(1-\theta_E)\log N}\right),
\]
which exceeds every power of $\log N$ for each fixed $\theta_E<1$. On
$\theta_E\le1-\theta'$ the support condition admits no $d\ge1$ and the
statement is vacuous, so what is established is the range
$\theta_E\in(1-\theta',1)$.
\end{theorem}

\begin{proof}
Identical to Theorem~\ref{thm:D} with $N^{1/2-\delta}$ replaced by
$N^{\theta_E}$: the residual's moduli are $md$ with $m<N^{1-\theta'}$,
so $d\le N^{\theta_E}/N^{1-\theta'}$, whence
$K/D=N^{\theta'}/N^{\theta_E-1+\theta'}=N^{1-\theta_E}$, and
Lemma~\ref{lem:extract} gives the ratio.

This asks $\theta_E>1-\theta'$. Below that the support bound
$N^{\theta_E-(1-\theta')}$ is under $1$, so no $d\ge1$ meets it and the
hypothesis on $b$ is empty --- in the same way
Note~\ref{note:nogoscope} records for Theorem~\ref{thm:D}. The
displayed exponent is therefore established on $\theta_E>1-\theta'$ and
the statement is vacuous below; for $\theta'>1/2$ the region left out
is $\theta_E<1/2$, beneath what Bombieri--Vinogradov already gives.
\end{proof}

Closing the gap would require $\theta_E=1$ exactly: equidistribution
of $\Lambda$ in progressions to moduli of size $N$ itself, where each
progression contains $O(1)$ terms and the statement carries no
information. So the route stays closed \emph{even granting the full
Elliott--Halberstam conjecture} --- it needs not a stronger level but
a different mechanism.

\begin{measurement}[the load-bearing lemma, checked as an identity]%
\label{meas:extract}
Lemma~\ref{lem:extract} carries the theorem, so it is checked
directly at $N=99\,999\,998$, $\theta'=0.56$,
$K=\lfloor N^{\theta'}\rfloor=30199$. For the single-divisor weights
$w_k=[\,d_0\mid k\,]$, i.e. $b=\delta_{d_0}$, the lemma is an identity
and is verified as one: over all $10\,262$ squarefree $d_0<K$ coprime
to $N$, the brute-force $B_w$ over $k<K$ and the factorised
$\mu(d_0)\varphi(d_0)^{-1}\rho_{d_0N}(K/d_0)$ agree to
$6.505\cdot10^{-18}$, the worst case being $d_0=33$.

The size behaves as the lemma predicts. The quantity
$|B_w|\varphi(d_0)$ is of order $1$ only when $K/d_0=O(1)$, i.e. only
when the weight's transform sits at the truncation point, and is
damped throughout the range Bombieri--Vinogradov admits: for these
parameters $N^{\theta'-1/2}=3$, and over $d_0\le3$ the largest value
attained is $0.004145$. The value is not a function of $K/d_0$ alone
--- $\rho_{d_0N}$ carries the condition $(j,d_0N)=1$, so it depends on
which small primes divide $d_0$. At $\lfloor K/d_0\rfloor=7$ the $184$
admissible $d_0$ produce exactly four values, split by $\gcd(d_0,15)$:
\[
\begin{array}{r|cccc}
 \gcd(d_0,15) & 1 & 3 & 5 & 15\\\hline
 |B_w|\,\varphi(d_0) & 0.250000 & 0.750000 & 0.500000 & 1.000000
\end{array}
\]
and at $\lfloor K/d_0\rfloor=13$ the $55$ admissible $d_0$ produce nine
values spanning $0.066667$ to $0.816667$, the largest at
the $d_0$ divisible by $15$ and by no other prime below $14$, where
$j=3$ and $j=5$ drop out of $\rho_{d_0N}$ and what is left is
$1-\tfrac1{10}-\tfrac1{12}$
(\texttt{audit\_extraction\_tradeoff.py}). The list stood here as six
values to $0.566667$ until the run behind it was found to be printing
only the first six of its own column, without saying so; the three
it dropped --- $0.650000$, $0.666667$, $0.816667$ --- were the three
largest. So the decay in $\lfloor K/d_0\rfloor$
is slower than this measurement was read as saying --- which does not
touch Theorem~\ref{thm:D}, whose range is $d_0\le3$, i.e.
$\lfloor K/d_0\rfloor\ge10066$, where the maximum is $0.004145$.
\end{measurement}

% =====================================================================
\section{Smooth weights:
\texorpdfstring{$\mu*\log^D=\Lambda_D$}{mu * log\^{}D = Lambda\_D}}%
\label{sec:Dpp}

Theorem~\ref{thm:D} closes the weights whose transform $b=\mu*w$ is
supported low enough for Bombieri--Vinogradov. That hypothesis
excludes the most natural family of all, $w_k=f(\log k)$ with $f$ a
polynomial, whose transform is spread over every scale. For those the
structure is explicit:
\[
  b\;=\;\mu*\log^D\;=\;\Lambda_D,
\]
the generalised von Mangoldt function. $\Lambda_D$ vanishes on
integers with more than $D$ prime factors, and on a squarefree
$u=p_1\cdots p_r$ with $r\le D$ it is the $r$-th finite difference of
$x^D$ at the points $\log p_i$; in particular
$\Lambda_D(u)=D!\,\log p_1\cdots\log p_D$ when $r=D$. Hence for
$f=x^D$ the complete part splits by the number of prime factors,
\begin{equation}\label{eq:CPD}
  \mathrm{CP}_D(N)=\sum_{r=1}^{D}\
    \sum_{\substack{u<N\ \text{squarefree}\\ \omega(u)=r}}
   \Lambda(N-u)\,\Lambda_D(u),
\end{equation}
the $r=1$ piece being a Goldbach-type sum and the $r\ge2$ pieces
Chen-type sums (prime plus $r$-almost-prime).

\begin{proposition}[polynomial weights]\label{prop:Dpp}
Let $w_k=f(\log k)$ with $f=\sum_{j\le D}c_jx^j$ real.
\begin{enumerate}
\item[(i)] If $D=0$ then $B_w\ll e^{-c\sqrt{\log K}}$ (Huang--Li's
Lemma~1): no extraction.
\item[(ii)] If $f$ is a single monomial $x^D$, $D\ge1$, then every
term of $\mathrm{CP}_D$ is nonnegative, since $\Lambda\ge0$ and
$\Lambda_D\ge0$; so $\mathrm{CP}_D\ge0$ and \emph{no cancellation is
available inside it}. Classical sieve upper bounds are expected to give
$\mathrm{CP}_D\ll N(\log N)^{D-1}$; that bound is quoted as standard,
is not proved here, and nothing below rests on it. A matching lower
bound is not
available: at $D=1$ the sum is exactly the binary Goldbach sum
$\sum_{p<N}\Lambda(N-p)\log p$, whose order is the assertion to be
proved. So $\mathrm{CP}_D$ is not known to be $o(N)$, and at $D=1$
deciding whether it is would settle Goldbach.
\item[(iii)] Cancellation therefore requires coefficients of opposite
sign, i.e. tuning $c_1$ against $c_2,\dots$; but the ratio that would
have to be matched involves the asymptotics of the $r=1$ pieces
$\sum_{p<N}\Lambda(N-p)\log^jp$, whose order is unknown for every
$j\ge1$ and which at $j=1$ \emph{is} the binary Goldbach sum. The
tuning is circular.
\end{enumerate}
Consequently no polynomial weight makes the complete part a
classically known object.
\end{proposition}

\begin{proof}
(i) is Huang--Li's Lemma~1 applied to $w\equiv1$. For (ii),
$\Lambda_D\ge0$ on the squarefree integers with $\omega\le D$ because
for $u=p_1\cdots p_r$ the value $\Lambda_D(u)$ is the $r$-fold mixed
difference of $x\mapsto x^D$ at $\log p_1<\cdots<\log p_r$, and an
$r$-fold difference of a function whose $r$-th derivative
$D(D-1)\cdots(D-r+1)x^{D-r}$ is nonnegative is itself nonnegative;
convexity alone would give this only for $r\le2$. Every summand of
\eqref{eq:CPD} is then a product of nonnegative factors, so no
cancellation is available and $\mathrm{CP}_D\ge0$. The upper bound
quoted in the statement is the classical sieve upper bound for
prime-plus-almost-prime sums applied to each $r$-piece; it is not
proved here and (ii) does not use it --- what closes (ii) is
nonnegativity together with the absence of a matching lower bound.
No lower bound of the same order is claimed:
the $r=1$ piece is a prime-plus-prime count, on which the sieve gives
none --- that is the parity obstruction of Note~\ref{note:bombieri} ---
and at $D=1$ it is the whole of $\mathrm{CP}_1$. For
(iii), a degree-$D$ tuning has $D$ free coefficients and the
constraint that must be solved is an identity between the leading
asymptotics of the $r$-pieces; the $r=1$ piece with $D=1$ is
$\sum_{p<N}\Lambda(N-p)\log p$, whose asymptotic is the assertion to
be proved.
\end{proof}

\begin{note}[the canonical tuning moves the wrong way]
Since $b=\mu*w$ we have $w=\mathbf 1*b$ and hence
$B(s)=W(s)/\zeta(s)$. That series has a \emph{double} pole at $s=1$
for every degree-$2$ $f$, so no such $f$ removes it and $b$ cannot be
made mean-zero; the most a degree-$2$ tuning can do is kill the
second-order term. With $f=x^2+cx$ one has $W=\zeta''-c\zeta'$, and
neither $\zeta''$ nor $\zeta'$ carries a $(s-1)^{-1}$ term.

Which way it moves is measured, and the title is not rhetorical. At
$N=10^6$ the untuned $\mathrm{CP}$ is $39.8793\,N$; the canonical
tuning $x^2+(2+2\gamma)x$ \emph{raises} it to $45.4201\,N$, a rise of
$13.89$ per cent, while $x^2-2\gamma x$ lowers it to $37.8516\,N$, by
$5.08$ per cent. The same ordering holds at $4\cdot10^6$ and
$1.6\cdot10^7$ (\texttt{audit\_polyweight.py}). Neither tuning buys
smallness, and the canonical one is the worse of the two: killing the
mean of $b$ is not what is wanted. What must be small is $b$'s
correlation with $\Lambda(N-\cdot)$, not its mean.
\end{note}

\begin{measurement}[the two pieces are the same order]\label{meas:poly}
Nonnegativity, and not a separation of scales, is what closes (ii).
At $N=10^6,\,4\cdot10^6,\,1.6\cdot10^7$ the two pieces of
$\mathrm{CP}_2$ are the same order:
\[
\begin{array}{r|ccc}
 N & 10^6 & 4\cdot10^6 & 1.6\cdot10^7\\\hline
 r=1:\ \sum_p\Lambda(N-p)\log^2p\,/\,N & 22.5145 & 25.0438 & 27.4579\\
 r=2:\ 2\sum_{pq}\Lambda(N-pq)\log p\log q\,/\,N
   & 17.3648 & 19.7926 & 22.2512\\
 \text{ratio } r_2/r_1 & 0.7713 & 0.7903 & 0.8104\\
 \mathrm{CP}_2/(N\log N) & 2.8866 & 2.9494 & 2.9967
\end{array}
\]
The ratio drifts toward $1$, not toward $0$ or $\infty$. Calibration:
the $D=1$ column reproduces the Goldbach sum
$\sum_p\Lambda(N-p)\log p=1.7565N,\,1.7633N,\,1.7614N$ against
$\SS(N)=1.76043$, which is the same at all three $N$ because
$N$ has the same prime support $\{2,5\}$ in each case
(\texttt{audit\_polyweight.py}).
\end{measurement}

% =====================================================================
\section{Both standard pairings have zero margin on
\texorpdfstring{$C(N)$}{C(N)}}\label{sec:circle}

Since the switch is closed over the part of its design space
Bombieri--Vinogradov evaluates, and of the remainder only the
polynomial weights are reached, in the weaker sense of
Proposition~\ref{prop:Dpp}, it is worth
recording the exact position of the other classical mechanism. Write
\[
  C(N)=\int_0^1S_\Lambda(\alpha)\,S_\mu(\alpha)\,e(-N\alpha)\,d\alpha,
\]
where
\[
  S_\Lambda(\alpha)=\sum_{n\le N}\Lambda(n)e(n\alpha),
  \qquad
  S_\mu(\alpha)=\sum_{m\le N}\mu(m)e(m\alpha).
\]

\begin{proposition}[zero margin]\label{prop:E}
The two pairings a circle-method estimate of this integral normally uses
--- $(2,2)$ and $(\infty,1)$ --- lie at or above the
trivial bound $|C(N)|\le\psi(N)\sim N$:
\begin{enumerate}
\item[(i)] $\|S_\Lambda\|_2\|S_\mu\|_2
   =\bigl(\sum_{n\le N}\Lambda(n)^2\bigr)^{1/2}
     \bigl(\sum_{m\le N}\mu^2(m)\bigr)^{1/2}
   \sim(6/\pi^2)^{1/2}\,N(\log N)^{1/2}$, exceeding the trivial bound
   by a factor $\asymp(\log N)^{1/2}$, which \emph{grows};
\item[(ii)] any bound of the shape
  $\sup_\alpha|S_\mu(\alpha)|\cdot\|S_\Lambda\|_1$ is at least
  $\|S_\mu\|_2\|S_\Lambda\|_1\gg N^{1/2}\cdot N^{1/2}=N$.
\end{enumerate}
\end{proposition}

\begin{proof}
(i) is Cauchy--Schwarz followed by
$\sum_{n\le N}\Lambda(n)^2\sim N\log N$ and
$\sum_{m\le N}\mu^2(m)\sim6N/\pi^2$. For (ii), the two factors are
bounded below for different reasons. Parseval gives
$\sup_\alpha|S_\mu|\ge\|S_\mu\|_2=(6N/\pi^2)^{1/2}$, an
identity-level constraint: no improvement in M\"obius
exponential-sum technology can lower $\sup_\alpha|S_\mu|$ below
$(6/\pi^2)^{1/2}N^{1/2}$. For the second factor Parseval runs the
\emph{wrong way}: it gives $\|S_\Lambda\|_1\le\|S_\Lambda\|_2
\ll(N\log N)^{1/2}$, an upper bound, while the trivial lower bound
$\|S_\Lambda\|_1\ge\bigl|\int_0^1S_\Lambda(\alpha)e(-n\alpha)
d\alpha\bigr|=\Lambda(n)$ reaches only $\log N$. The bound
$\|S_\Lambda\|_1\gg N^{1/2}$ that (ii) uses is a theorem of Vaughan
\cite{Vau88}.
\end{proof}

The classical uniform bound $S_\mu(\alpha)\ll_AN(\log N)^{-A}$, due to
Davenport and quoted here as standard rather than proved, is
therefore useless here: it is a saving over $N$, whereas the pairing
needs a bound at the scale $N^{1/2}$, which Parseval forbids. This is
the parity obstruction in circle-method language --- \emph{the binary
problem sits exactly at the trivial bound with no margin} --- and it
is why the method that settles the ternary problem cannot be pushed to
the binary one by any sharpening of the M\"obius input.

\begin{measurement}[the constants]\label{meas:circle}
Computed by exact FFT on a $4N$-point grid at
$N=2^{14},\dots,2^{20}$:
\[
\begin{array}{r|cccc}
 N & 2^{14} & 2^{16} & 2^{18} & 2^{20}\\\hline
 \|S_\mu\|_2/\sqrt N & 0.7798 & 0.7797 & 0.7797 & 0.7797\\
 \sup|S_\mu|/\sqrt N & 3.058 & 2.742 & 2.853 & 2.801\\
 \|S_\Lambda\|_1/\sqrt N & 1.946 & 2.084 & 2.219 & 2.346\\
 \text{(i) bound}/N & 2.297 & 2.473 & 2.639 & 2.795\\
 N/\bigl(\sup|S_\mu|\,\|S_\Lambda\|_1\bigr) & 0.168 & 0.175 & 0.158
  & 0.152
\end{array}
\]
The first row reproduces $\sqrt{6/\pi^2}=0.779697$ to four significant
figures --- the entries before rounding are $0.779764$, $0.779725$,
$0.779686$, $0.779699$ --- confirming
the computation; the last row --- the margin that route~(ii) would
need to exceed $1$ --- sits near $0.16$ and falls over the last three
abscissae. Route~(i) diverges from the trivial bound like
$(\log N)^{1/2}$, as predicted. For reference the object itself is
small: $C(N)/N=-0.0105,\,0.0001,\,0.0059,\,0.0032$ at these $N$
(\texttt{audit\_circle\_margin.py}).
\end{measurement}

% =====================================================================
\section{Summary}

\begin{itemize}
\item The design space of divisor-switch weights is described exactly
by six finite rearrangements ---
Propositions~\ref{prop:dilate}, \ref{prop:posweights},
\ref{prop:flatsum}, \ref{prop:untrunc}, \ref{prop:layers} and
\ref{prop:combined} --- with
Proposition~\ref{prop:direct} a consequence of them rather than one of
them: the per-modulus
discrepancy is a dilate, the signs $\mu(k)$ cancel on the $\log k$
branch, and the untruncated object is the Goldbach count itself.

\item Theorem~\ref{thm:D}: no weight whose transform is
Bombieri--Vinogradov-accessible extracts $C(N)$ by the evaluation this
route offers, and this holds even
when the bound on the switch identity's third term $T_w$ --- the object
the whole problem is about --- is granted for free. What is closed is
the triangle-inequality evaluation; an evaluation through sign
correlations among the identity's three terms is not. The loss is
$\exp(c_1\sqrt{(\frac12+\delta)\log N})$, and the $1/2$ is the exponent of the
square-root barrier.

\item Theorem~\ref{thm:Dprime}: the separation survives
$EH(N^{\theta_E})$ for every fixed $\theta_E<1$ at which the
hypothesis is non-empty, that is on $\theta_E\in(1-\theta',1)$.
Closing it would require $\theta_E=1$, where the statement is vacuous.

\item Proposition~\ref{prop:Dpp}: polynomial weights, which lie
outside the hypothesis of Theorem~\ref{thm:D}, are met by
nonnegativity of $\Lambda_D$ and by circularity of the tuning ---
which leaves their complete part not known to be $o(N)$ rather
than shown to be large. That is a statement about what is not
known and not an impossibility, and the weights that are neither
reached by the hypothesis nor polynomial are closed nowhere here.

\item Proposition~\ref{prop:E}: the two pairings a circle-method
estimate of $C(N)$ normally uses have zero margin, with Parseval
supplying the floor. A bilinear treatment of the minor arcs is a
different mechanism and is not addressed.

\item Net progress toward Goldbach: zero. What is closed is a design
space, so that it need not be re-explored.
\end{itemize}

% =====================================================================
\begin{thebibliography}{99}

\bibitem{Bom76} E.~Bombieri, \emph{The asymptotic sieve},
Rend. Accad. Naz. XL (5) \textbf{1(2)} (1975/76), 243--269 (1977).

\bibitem{GY} D.~Goldston and C.~Y{\i}ld{\i}r{\i}m,
\emph{Higher correlations of divisor sums related to primes~I},
Integers \textbf{3} (2003), A5.

\bibitem{HL} J.-J.~Huang and H.~Li, \emph{On the connection between
the Goldbach conjecture and the Elliott--Halberstam conjecture},
arXiv:2005.03811 [math.NT]; v1 (May 2020), v2 (August 2022). Also
published in a Springer volume,
\texttt{doi:10.1007/978-3-030-67996-5\_17}. References here are to the
arXiv version, which we have read, and not to the published version,
which we have not.

\bibitem{Li23} J.~D.~Lichtman (with an appendix by S.~Drappeau),
\emph{Primes in arithmetic progressions to large moduli, and Goldbach
beyond the square-root barrier}, arXiv:2309.08522 [math.NT], 2023.

\bibitem{MoV} H.~L.~Montgomery and R.~C.~Vaughan,
\emph{Multiplicative number theory~I: classical theory}, Cambridge
Studies in Advanced Mathematics \textbf{97}, Cambridge University
Press, 2007.

\bibitem{P1} Y.~Lee, \emph{The Huang--Li reduction of Goldbach to
Elliott--Halberstam does not decompose, with an unconditional bound for
a M\"obius-weighted correlation sum in fixed residue classes}, 2026.
Companion paper, in the directory \texttt{P1-mobius-fixed-class/} of
the repository named in Appendix~\ref{sec:repro}, with its own code and
output. The switch identity \eqref{eq:extract} is its Theorem~C and is
not reproved here.

\bibitem{Vau88} R.~C.~Vaughan, \emph{The $L^1$ mean of exponential
sums over primes}, Bull. London Math. Soc. \textbf{20} (1988),
121--123. Not consulted in the original; the statement used here is
taken from its review and from two independent restatements.

\end{thebibliography}

% =====================================================================
\section*{Disclosure statement}

No potential competing interest, financial or non-financial, was
reported by the author. No funding was received for this work. An AI
assistant was used in preparing it; what tool, how, and why is set out
in the appendix, together with what was machine-verified, what was
checked by computation, and what was neither.

\section*{Data availability statement}

The code behind every figure printed here, the output it produced, and
the Lean~4 development are openly available at
\texttt{https://github.com/soke91/arithmetic-convolution-fields}, in the
directory \texttt{P2-no-go-divisor-switch/}. Two percentages quoted in
Measurement~\ref{meas:flatsum} are read off two printed columns of a
cited result file rather than produced by a run, and that passage says
which columns and at which $N$. No data derived from human or animal
subjects were used; every quantity here is computed from the integers.

% =====================================================================
\appendix
\section{Use of an AI assistant, and what checked what}

This work was done in close collaboration with an AI assistant, which
carried much of the derivation and all of the computation. The author
takes full responsibility for every claim made here.

\emph{Which tool, and how.} The manuscript in its present form was
prepared with Anthropic's Claude, model Opus~5, run through the Claude
Code command-line interface. It was used to derive and check the
identities, to write and run every script in the accompanying code, to
write the Lean development, and to draft and revise this text. Earlier
stages used earlier models of the same family; the repository does not
record which, and nothing is relied on from those stages that has not
since been rederived and rerun.

What follows says what was checked and by what, rather than asserting
that it was checked.

\emph{Machine-verified.} The finite rearrangements that lay out the
design space are formalised in Lean~4 against Mathlib, in the file
\texttt{lean/GoldbachLean/Layers.lean} of the code packet: the exchange
of the double sum along its inner variable that
Proposition~\ref{prop:layers} performs, the nonnegativity of the layers
it produces, and the sense in which the two ends of
Proposition~\ref{prop:flatsum} are one expression. Each depends on
\texttt{propext}, \texttt{Classical.choice} and \texttt{Quot.sound} and
on nothing else; the axiom report is in the packet. The Lean statements
take $\Lambda$, $\mu$ and $\log$ as arbitrary functions, so what is
verified is that these identities use no arithmetic property of them ---
which is what makes them rearrangements rather than estimates.

\textbf{The no-go itself is not machine-verified, and neither theorem of
this paper has had a reading by a subject expert.}
Theorem~\ref{thm:D} rests on Bombieri--Vinogradov, which Mathlib does
not carry, so no part of it is formalised. What is formalised is the
scaffolding, not the conclusion.

\emph{Checked by computation.} Every figure quoted is produced by a
named run in the accompanying code, with the two exceptions named in
Measurement~\ref{meas:flatsum}, which are read off two printed columns
of a cited file. A check in the build refuses a printed figure that no
cited run produces. Forty-three rules carry a verdict across the
fourteen cited runs, and every one is reported with it, including the
ten the runs did not meet.

\emph{Cross-read.} The proofs and the computations were additionally
read by separate AI sessions working without sight of each other's
findings, and several statements here are corrections those readings
produced --- among them the reference against which the cross-$k$ gain
is compared, which an earlier printing took to be a count of moduli. It
is not a substitute for either of the checks named as absent above.

\section{Reproducibility}\label{sec:repro}

The code and the output it produced are at

\begin{center}
\texttt{https://github.com/soke91/arithmetic-convolution-fields}
\end{center}

\noindent in the directory \texttt{P2-no-go-divisor-switch/}. Each entry
below is run as \texttt{python code/<script>.py} from that directory and
writes the file of the same name in \texttt{results/}. The figures here
were produced under Python~3.12.10 and NumPy~2.2.5. That address is
mutable, so a number checked against the branch head may have been
recomputed since; the release named there pins the state these figures
were read from.

Two files are in that directory without being named below:
\texttt{lab\_layer\_tail.py}, whose output
\texttt{lab\_combined\_modulus.py} reads, and \texttt{indep.py}, which
\texttt{c02\_flatsum\_k1.py} imports. Run \texttt{lab\_layer\_tail}
before \texttt{lab\_combined\_modulus}; the rest are independent of each
other.

\emph{A nonzero exit status is not a failure.} These runs carry
preregistered rules, and a run whose own rule was refuted says so in its
exit code as well as in its output. Eight of the fourteen runs named
below record such a refutation, and all eight are named in
\S\ref{sec:notclaimed}.

\begin{center}
\small
\begin{tabular}{llp{3.8cm}}
\hline
Script & Result file & Supports\\
\hline
\texttt{lab\_dilate\_identity.py} & \texttt{lab\_dilate\_identity.txt}
 & Proposition~\ref{prop:dilate}\\
\texttt{lab\_positive\_weights.py} & \texttt{lab\_positive\_weights.txt}
 & Proposition~\ref{prop:posweights}\\
\texttt{lab\_weight\_gap.py} & \texttt{lab\_weight\_gap.txt}
 & Measurement~\ref{meas:flatsum}, $k\ge2$ column\\
\texttt{c02\_flatsum\_k1.py} & \texttt{c02\_flatsum\_k1.txt}
 & Measurement~\ref{meas:flatsum}, $k\ge1$ column\\
\texttt{lab\_direct\_identity.py} & \texttt{lab\_direct\_identity.txt}
 & Proposition~\ref{prop:untrunc}\\
\texttt{lab\_layer\_decomposition.py} &
 \texttt{lab\_layer\_decomposition.txt} & Proposition~\ref{prop:layers}\\
\texttt{lab\_combined\_modulus.py} & \texttt{lab\_combined\_modulus.txt}
 & Proposition~\ref{prop:combined}\\
\texttt{lab\_direct\_route.py} & \texttt{lab\_direct\_route.txt}
 & Measurement~\ref{meas:direct}, residual and parenthetical list\\
\texttt{a1\_bh\_vs\_b.py} & \texttt{a1\_bh\_vs\_b.txt}
 & Measurement~\ref{meas:direct}, the two lists\\
\texttt{audit\_extraction\_tradeoff.py} &
 \texttt{audit\_extraction\_tradeoff.txt} & Measurement~\ref{meas:extract}\\
\texttt{audit\_polyweight.py} & \texttt{audit\_polyweight.txt}
 & Measurement~\ref{meas:poly}\\
\texttt{audit\_circle\_margin.py} & \texttt{audit\_circle\_margin.txt}
 & Measurement~\ref{meas:circle}\\
\texttt{audit\_nogo\_scope.py} & \texttt{audit\_nogo\_scope.txt}
 & Note~\ref{note:nogoscope}\\
\texttt{audit\_squarefull\_admitted.py} &
 \texttt{audit\_squarefull\_admitted.txt} & the note after
 Proposition~\ref{prop:flatsum}\\
\hline
\end{tabular}
\end{center}

No computation is used as a premise of any proof. Six of the seven
statements of Section~\ref{sec:identities} are finite rearrangements
and are verified as such to machine precision;
Proposition~\ref{prop:direct} is the seventh and is not one of them ---
it carries an $O_A(N(\log N)^{-A})$ --- so what is verified there is
its identity part and not the error term. The measurements record
sizes over the ranges stated in each statement.

Eight of the fourteen runs record a pre-registered rule their own
output did not meet, and this paper reported none of them by name.
They are:
\texttt{audit\_extraction\_tradeoff} (two rules, one of them the damping
claim itself --- above $\lfloor K/d_0\rfloor=1000$ the maximum is
$0.062792$ against a registered cap of $0.05$; the range
Theorem~\ref{thm:D} uses is $\lfloor K/d_0\rfloor\ge10066$, where the
maximum is $0.004145$, so the theorem is untouched and the registered
cap was set too tight);
\texttt{lab\_combined\_modulus} (two rules), whose control with the sign
of $\mu(d)$ randomised has a \emph{smaller} median deficit than $\mu$
at three of its four $N$, including the two largest --- $0.0454$
against $0.0547$ and $0.0342$ against $0.0733$ --- and a smaller
minimum at every one, so that measurement does not separate $\mu$ from
a sign pattern at all in the range it covers, and what makes the far
pairs negligible is their size rather than their signs;
\texttt{lab\_direct\_route} (two rules), residual $0.2076$ against a
registered $0.20$; its randomised-sign control is inverted in the same
way, the coin residuals having median $0.0157$ against $\mu$'s
$0.2076$ --- smaller by a factor $13.2$ --- and that run's own
diagnostic records that the null as designed tests nothing, since for
a coin $\sum(\log k)H_\varepsilon$ and $\tilde r-\SS N$ are separately
near zero and their difference is small for a reason unconnected with
the cancellation;
\texttt{lab\_positive\_weights}, top-decile share near $35\%$ against a
registered $40\%$;
\texttt{lab\_direct\_identity} and \texttt{lab\_layer\_decomposition},
each on a rule about how much mass sits below the truncation;
\texttt{audit\_nogo\_scope}, whose rule TD3 asked one hypothesis to
cover both $b=\mu*1$ and $b=\mu*\log$ and it covers neither, which is
what Note~\ref{note:nogoscope} says in words;
and \texttt{audit\_squarefull\_admitted}, whose V3 asked the note after
Proposition~\ref{prop:flatsum} to reproduce its own printed
three-decimal figures and one of them was a double rounding, corrected
above to what that run reads.
Of the eight, \texttt{lab\_direct\_route}'s Y1 is the one that bears on
a statement: it is the convergence
$B_{\log}(K)+\SS(N)\to0$ that Proposition~\ref{prop:direct} uses, and
what its failure records is that the run does not exhibit the rate at
its own $N$, not that the limit is false. The other seven change no
statement above --- each was checked --- and the reason for recording
all of them here is that a reader who does not open the files would
otherwise take the runs for clean.

\end{document}
